\begin{established}[The split refinement of the marked core]\label{hyp:split-section-refinement}
In the local construction write \(p=\ell\), \(M=g_0g_1\) for
the original full level and \(I=P\mathcal O_E\), where \(P\)
is the chosen prime above \(p\) in the marked imaginary
quadratic field. If \(J\) denotes the old half, then the new level data are
\[
 (\text{level},\text{chosen half},\text{auxiliary level},\text{field})
                       =(Mp,JI,1,F),
\]
with the same coefficient field \(F\). For each \(\sigma\in\operatorname{Aut}(\mathbb C/F)\), the marked CM curve and its pullback are isomorphic by an isomorphism respecting every torsion label at level \(Mp\) and at each divisor of \(Mp\).
\end{established}
\begin{proof}
We have \(I+\overline I=\mathcal O_E\) and
\(I\overline I=(p)\), by splitting. Use at each lift the
polarization and moduli problem of the core \(J\) just constructed.
At each \(j\) refine the split level by \(I,\overline I\).
Writing \(J\) for the old chosen half, the updated parameters are
\[
 (\text{integer},\text{chosen half},\text{auxiliary level})=(Mp,JI,1).
\] The two complementary halves exist by the displayed
factorization and coprimality. Use the moduli space with this
refined marking and the map forgetting its new level.
The coefficient field stays fixed. The field contains the specified roots of the coprime orders \(M\) and \(p\); their product has order
\[
 \operatorname{ord}(\zeta_M\zeta_p)=\operatorname{lcm}(M,p)=Mp .
\]

Write \(F_{\rm rec}\) for the embedded record field and \(F\) for
the cover coefficient field. The pin on the coefficient field
sends the generators of \(F_{\rm rec}\) to their prescribed complex
values in \(F\), so \(F_{\rm rec}\subseteq F\).
For the marked CM lattice \(\Lambda\) the field conditions read
\[
 k\big(j(\Lambda), W(P):P\in (Mp)^{-1}\Lambda/\Lambda\big)\subseteq F .
\]
Here \(k\) denotes the marked imaginary quadratic field and \(W\) its Weber function. 

For each automorphism
\(\sigma\) of \(\C\) fixing \(F\), these values are fixed. By the main theorem of complex multiplication in its ray-class form \cite{cm-reciprocity}, the Weber values at level \(Mp\) lie in its ray field, contained in \(F\) by Lemma~\ref{lem:marked-original-geometric-field}. The equality of \(j\)-invariants gives an isomorphism to the pullback. Since \(F\) contains the CM field, it respects the specified CM action. Let \(P_0\) generate the free rank-one \(\mathcal O_k/(Mp)\)-module of torsion. Write \(f\) for the isomorphism to the pullback and \(W\) for its Weber quotient. Since \(W\) is the quotient by the CM units, the equality of the two readings gives
\[
 W(f(P_0))=W(P_0^\sigma)
 \quad\Longrightarrow\quad f(P_0)=uP_0^\sigma
                  \quad(u\in\mathcal O_k^\times).
\]

The map \(u^{-1}f\) now matches \(aP_0\) for every \(a\in\mathcal O_k/(Mp)\) by CM-linearity. This isomorphism may depend on \(\sigma\). For any \(m\mid Mp\) the
same isomorphism works at level \(m\), since the point with
label \(a/m\) has label \(((Mp/m)a)/(Mp)\) at the original
level.
\label{cor:split-level-sections-1}
\end{proof}

\begin{proposition}[A fixed split refinement for the short theta rules]\label{prop:split-level-sections}
For the split prime and the marked core \(J\) of Lemma~\ref{hyp:split-section-refinement},
one can choose the same cover, coefficient field and weight-zero
determinant character for every number \(n\) of torsion arguments.
For \(n\ge r_j=p_j+q_j\) the Fourier coefficient is obtained from
the basis test below, with the inverse character under a change of
marking. If \(h_n\) denotes the coefficient attached to \(n\) torsion arguments, the rank obstruction gives
\[
 n<r_j\quad\Longrightarrow\quad h_n=0.
\]
\end{proposition}

\begin{proof}
For the component comparison the refined auxiliary level is
one, the rank is positive, and the marked lattice supplies a
rank-one definite datum for the same host. In indefinite
signature the rank is at least three. Under these hypotheses
the CM comparison for the unitary period spaces says that the
pullback isomorphism fixing the torsion labels and the host
root induces the identity permutation of the period
components. Consequently an algebraic constant in the
function field of a component is fixed by every
\(\sigma\in\operatorname{Aut}(\C/F)\) under its generic
embedding. Indeed, if an algebraic constant \(a\) had degree greater than one, separability and extension of embeddings to \(\mathbb C\) would give
\[
 \sigma\in\operatorname{Aut}(\mathbb C/F),\qquad \sigma(a)\ne a,
\]
contrary to the asserted invariance. Thus \(F\) is relatively algebraically closed in
each component function field. The component idempotents and
the prescribed values in \(F\) now give the componentwise
weight-zero algebraic rules.
\label{cor:split-level-sections-2}

\proofstep{The determinant character.}
Here one must put a nontrivial character on that refinement.
Let \(\mathbb A_{E,f}^1\) be the norm-one subgroup of the
finite ideles. Hilbert 90 gives
\[
 1\longrightarrow\mathbb A_{F,f}^\times
 \longrightarrow\mathbb A_{E,f}^\times
 \xrightarrow{\,h:y\mapsto y/\bar y\,}\mathbb A_{E,f}^1
 \longrightarrow1 .
\]
The finite character \(d_p\) read from \(\delta_p\) is trivial on
the left-hand subgroup. In taking the finite part its
archimedean kernel is also trivial, by divisibility of the
connected complex idele factors and finite order of the
character. Hence \(\xi_p(h(y))=d_p(y)\) is independent of the
lift \(y\). A field-valued complex character on this
subgroup extends as a \(\C^\times\)-valued character of the
abelian group, by divisibility of \(\C^\times\). Only its values
on norm-one determinants will be used. A rational norm-one
determinant has a rational lift by Hilbert 90, so \(\xi_p\)
is trivial there.

At principal level the local congruence
version provides an integral lift congruent to one at \(p\).
At the odd split special prime the componentwise involution interchanges factors, so the lift \(z=(u,1)\) gives
\[
 z/\bar z=(u,1)/(1,u)=(u,u^{-1}).
\] The tame
character therefore gives one on the required level kernel,
and its nontrivial residue value at \(p\) survives on the
determinant quotient. Its values on norm-one determinants are
the original finite-order values in the coefficient field.
The function \(g\mapsto\xi_p(\det g)\) is constant on the
corresponding marked components, by the rational and level
invariance. The preceding component comparison gives a
weight-zero rule with these values, whose same positive power
is one on every component.
\label{cor:split-level-sections-3}

\proofstep{The finite Fourier coefficient.}
The finite Fourier expansion must retain the character just obtained on the determinant.

Write \(r_j=p_j+q_j\) and consider \(n\) torsion arguments on
this cover. For \(n\ge r_j\) fix a basis \(y\) of its second
split half module of rank \(r_j\). On an \(n\)-tuple \(x\)
whose first \(r_j\) entries form a basis \(g y\), set
\[
              h_n(x)=\psi(g^\dagger)^{-1}.
\]
Here \(\psi\) is the character on automorphisms of the first
half induced by the nontrivial determinant character, and
\(g^\dagger\) is the adjoint between the two dual halves.
Set \(h_n=0\) elsewhere. Over the finite residue quotient these
modules are free and a spanning \(r_j\)-tuple is a basis,
as can be checked at every local factor by its determinant.
The basis change \(g\) is unique. Thus \(h_n(y,0)=1\). The inverse character under change of marking follows from reversal of composition by adjunction:
\[
 (g'g)^\dagger=g^\dagger(g')^\dagger .
\]

Take the finite Fourier
sum of \(h_n\) against the additive character of the level
pairing in the \(n\) arguments. Its coefficients lie in the
field just specified. This is the finite coefficient of the
short theta rule at this order. When \(n<r_j\) its spanning
condition has empty hypothesis and we take the zero coefficient.
The weight-zero twist section, the underlying cover and its
coefficient field were fixed before \(n\). Only the coefficient
of the theta rule changes. Finite Fourier evaluation and the
harmonic chart transformation compute these short rules on the
same cover at all orders, with the same determinant twist.
\label{cor:split-level-sections-4}
\end{proof}
